
\question[3]
Eine Audiodatei wurde ursprünglich mit einer Samplingrate von 44\,100\,Hz aufgenommen
und soll nun mit 11\,025\,Hz gespeichert werden. Beschreibe, was bei dieser Reduktion 
der Samplingrate mit dem Signal passiert und welchen Nachteil dies mit sich bringen kann.

\begin{solutionbox}{5cm}
Bei der Aufnahme wird die Schallwelle in festen Zeitabständen abgetastet (gesampelt). 
Eine geringere Samplingrate bedeutet, dass pro Sekunde weniger Abtastwerte gespeichert werden
 -- hier nur noch ein Viertel der ursprünglichen Menge. 
Dadurch verringert sich die Dateigröße entsprechend.
Nachteilig ist, dass hohe Frequenzanteile des Signals nicht mehr korrekt erfasst werden
können wodurch der Klang weniger originalgetreu wirkt. 
\end{solutionbox}